\documentclass[10pt,letterpaper]{article}
\usepackage[utf8]{inputenc}
\usepackage[margin=0.6in]{geometry}
\usepackage{titlesec}
\usepackage{enumitem}
\usepackage{hyperref}
\usepackage{xcolor}
\usepackage{fontawesome5}

% Color Configurations
\definecolor{primary}{RGB}{30, 41, 59}   % Slate Blue
\definecolor{secondary}{RGB}{71, 85, 105} % Soft Charcoal
\definecolor{accent}{RGB}{2, 132, 199}   % Deep Sky Blue

\hypersetup{
    colorlinks=true,
    linkcolor=accent,
    urlcolor=accent
}

% Typography Configurations
\urlstyle{same}
\pagestyle{empty}

% Section Headers Formatting
\titleformat{\section}
  {\large\bfseries\color{primary}}
  {}{0em}
  {}[\color{secondary}\titlerule\vspace{0.4em}]

\titlespacing*{\section}{0pt}{1.2em}{0.4em}

% Custom List Formats
\setlist[itemize]{leftmargin=1.5em, noitemsep, topsep=0.2em, parsep=0.1em}

\begin{document}

% Add this line to the preamble of your document (before \begin{document}):
% \usepackage{fontawesome5}

% ==================== HEADER ====================
\begin{center}
    {\LARGE\bfseries\color{primary} Satyam Awasthi} \\ \vspace{0.3em}
    {\color{secondary} Researcher \& Software Engineer} \\ \vspace{0.4em}
    \small 
    \faEnvelope~\href{mailto:satyam.aw@yahoo.com}{satyam.aw@yahoo.com} \ | \ 
    \faGlobe~\href{https://satyam-aw.github.io}{satyam-aw.github.io} \ | \ 
    \faGithub~\href{https://github.com/satyam-aw}{satyam-aw} \ | \ 
    \faOrcid~\href{https://orcid.org/0000-0002-6095-4303}{0000-0002-6095-4303} \ | \ 
    \faLinkedin~\href{https://linkedin.com/in/satyam-awasthi}{LinkedIn}  \\ \vspace{6pt}
    % Kanpur, IN \,$\cdot$\, 
    % UCSB \& IIT Kharagpur Alum
\end{center}


\vspace{0.5em}

% ==================== RESEARCH SUMMARY ====================
\section*{Professional Summary}
Computer Science researcher specializing in \textbf{Autonomous Decision-Making}, \textbf{Spatial AI}, and \textbf{Embodied AI} within \textbf{Cyber-Physical Systems}, with a core focus on engineering verifiably safe, constraint-bound execution loops.

% ==================== EDUCATION ====================
\section*{Education}
\noindent
\textbf{University of California, Santa Barbara (UCSB)} \hfill Santa Barbara, CA \\
\textit{M.S. in Computer Science (GPA: 3.91/4.00)} \hfill Sep 2021 – Jun 2023 \\
Graduate Research focused on Spatial AI, AR/VR, and HCI at the Four Eyes Lab. \\
Advisor: \href{https://scholar.google.com/citations?user=008lr2cAAAAJ&hl=en}{\color{accent}Prof. Tobias Höllerer} \,$\cdot$\, Co-Advisor: \href{https://scholar.google.com/citations?user=dK-0kG4AAAAJ&hl=en}{\color{accent}Prof. Michael Beyeler} \\

\noindent
\textbf{Indian Institute of Technology, Kharagpur (IIT Kharagpur)} \hfill Kharagpur, India \\
\textit{B.Tech. in Electrical Engineering, Minor in Computer Science (GPA: 8.98/10.0)} \hfill Jul 2016 – Jul 2020 \\
Graduated with First Class Honors.

% ==================== SKILLS ====================
\section*{Skills}
\noindent \textbf{Programming Languages:} Python, C++, Kotlin, Java, MATLAB \\
\noindent \textbf{Frameworks \& Core Tools:} Android/iOS Native SDKs, LangChain, Pydantic (JSON Schema), OpenCV, ROS (Robot Operating System), PyTorch \\
\noindent \textbf{Fields of Interest:} Autonomous Decision-Making, Embodied AI, Spatial AI, Multi-Modal Sensor Fusion, Cyber-Physical Systems (CPS) 

% ==================== PUBLICATIONS ====================
\section*{Publications}

\subsection*{Refereed Conference Papers}

\noindent \textbf{SmartWatch for Wall Writing: Real-time Transcription of Wall Writing from Inertial Sensing} \\
Snigdha Das, \textbf{Satyam Awasthi}, Abdul Shamnar P, Pradipta De, Sandip Chakraborty, and Bivas Mitra. \\
\textit{Proceedings of the 14th International Conference on COMmunication Systems \& NETworkS (COMSNETS)}, 2022. \\
\href{https://doi.org/10.1109/COMSNETS53615.2022.9668531}{\color{accent}\footnotesize DOI: 10.1109/COMSNETS53615.2022.9668531} \\
Engineered an adaptive coordinate transformation pipeline using gravitational alignment and PUPD micro-gesture filtering to resolve drift-free, 2D kinematic profiles of wall-writing from wearable inertial sensors.

\medskip

\subsection*{Refereed Poster Papers}

\noindent \textbf{Eye Tracking Performance in Mobile Mixed Reality} \\
\textbf{Satyam Awasthi}, Vivian Ross, Sydney Lim, Michael Beyeler, and Tobias Höllerer. \\
\textit{IEEE Conference on Virtual Reality and 3D User Interfaces Abstracts and Workshops (VRW)}, 2024. \\
\href{https://doi.org/10.1109/VRW62533.2024.00321}{\color{accent}\footnotesize DOI: 10.1109/VRW62533.2024.00321} \\
Expanded the tracking framework to benchmark low-latency performance profiles and gaze errors via a robust 3-device user study across diverse mixed-reality locomotion layouts. Supported by NSF Grant \#2211784.

\smallskip
\smallskip

\noindent \textbf{EyeTTS: Evaluating and Calibrating Eye Tracking for Mixed-Reality Locomotion} \\
\textbf{Satyam Awasthi}, Vivian Ross, Michael Beyeler, and Tobias Höllerer. \\
\textit{IEEE International Symposium on Mixed and Augmented Reality Adjunct (ISMAR-Adjunct)}, 2023. \\
\href{https://doi.org/10.1109/ISMAR-Adjunct60411.2023.00104}{\color{accent}\footnotesize DOI: 10.1109/ISMAR-Adjunct60411.2023.00104} \\
Developed a novel \href{https://github.com/satyam-aw/Eye-Tracking-Test-Suite_EyeTTS_Data-Module}{eye-tracking test suite (EyeTTS)} and engineered unique software calibration pipelines optimized to correct tracking drift during active human locomotion; validated via a 2-device user study. Supported by NSF Grant \#2211784.

\newpage
% ==================== EXPERIENCE ====================
\section*{Experience}

\noindent
\textbf{Intuit Inc} \hfill Mountain View, CA \,$\cdot$\, Mar 2024 – Jan 2026 \\
\textit{Software Engineer 2} 
\begin{itemize}
    \item Built \href{https://satyam-aw.github.io/projects/intuit_generative_ui/}{dual-stage} \textbf{agentic decision pipelines} using tool-schema routing for real-time intent classification and sequential state tracking.
    \item Developed a decoupled \textbf{generative UI engine} leveraging RAG and structured LLMs to optimize deterministic component layouts.
\end{itemize}

\noindent
\textbf{Yahoo Inc} \hfill Mountain View, CA \,$\cdot$\, Apr 2023 – Nov 2023 \\
\textit{Software Dev Engineer II (IC3)} 
\begin{itemize}
    \item Developed \textbf{client-side rule engines} for asynchronous receipt payloads, executing deterministic filtering for \textbf{real-time decision support}.
    \item Designed a dynamic, state-driven carousel UI framework optimized for client-side rendering efficiency.
\end{itemize}

\noindent
\textbf{Coupang Global LLC} \hfill Mountain View, CA \,$\cdot$\, Jun 2022 – Sep 2022 \\
\textit{Software Engineering Intern} 
\begin{itemize}
    \item Engineered \textbf{context-aware heuristics} to prune multi-dimensional search spaces based on active user state vectors.
    \item Designed \textbf{Content-Based Filtering (CBF)} layers using \textbf{cosine similarity} to rank high-dimensional product features.
\end{itemize}

\noindent
\textbf{Samsung Research Institute} \hfill Noida, India \,$\cdot$\, Sep 2020 – Sep 2021 \\
\textit{Software Engineer (CL 2)} 
\begin{itemize}
    \item Architected a \textbf{continuous on-device perception pipeline} ingesting multi-modal media streams for automated ML classification.
    \item Designed a \textbf{cross-profile data brokering engine} leveraging multi-user storage isolation to route assets from \textbf{User 1} to sandboxed \textbf{User 0}.
\end{itemize}


% ==================== PROJECTS ====================
\section*{Projects}

\noindent
\textbf{\href{https://satyam-aw.github.io/projects/jitternot/}{JitterNot}} \hfill Jan 2022 – Mar 2022 \\
\textit{LSTM-Based Model Predictive Controller for Adaptive Video Streaming}  \hfill UCSB
\begin{itemize}
    \item Developed an LSTM network for real-time throughput prediction, transforming the closed-loop control system from MISO to SISO.
    \item Implemented an MPC controller leveraging Receding Horizon Control (RHC) principles to dynamically calculate optimal video resolution.
\end{itemize}

\noindent
\textbf{\href{https://satyam-aw.github.io/projects/autonomous_ground_vehicle_and_swarm_robotics/}{Autonomous Ground Vehicle \& Swarm Robotics}} \hfill Jan 2016 – Dec 2018  \\
\textit{End-to-End Perception, Planning, and Multi-Agent Navigation Stack}  \hfill IIT Kharagpur
\begin{itemize}
    \item Built an autonomous navigation stack in ROS, fusing heterogeneous sensor streams (LiDAR, IMU, GPS) via an Extended Kalman Filter (EKF).
    \item Implemented visual obstacle segmentation networks (U-Net) along with global ($A^*$/Dijkstra) and local (DWA) path planners.
    \item Researched decentralized swarm coordination using peer-to-peer relative pose estimation (AprilTags) over an ad-hoc XBee mesh network.
\end{itemize}


\noindent
\textbf{\href{https://satyam-aw.github.io/projects/in_motion_har_deepconvlstm/}{In Motion — HAR using DeepConvLSTM}}  \hfill Nov 2019 – Nov 2019 \\
\textit{Distributed Mobile Computing System for Human Activity Recognition} \hfill IIT Kharagpur 
\begin{itemize}
    \item Designed a client-server telemetry pipeline streaming high-frequency IMU sensor arrays for rolling window inference.
    \item Implemented a DeepConvLSTM core architecture using TimeDistributed CNN layers to model spatial-temporal features efficiently.
\end{itemize}

\noindent
\textbf{\href{https://satyam-aw.github.io/projects/cohesive_ar/}{CohesiveAR}}  \hfill Apr 2022 – Jun 2022 \\
\textit{Augmented Reality Pipeline for Texture Extraction and Homography Mapping} \hfill UCSB 
\begin{itemize}
    \item Developed a spatial perception pipeline using Google ARCore and OpenCV to map and project dynamic 3D scene geometry into unified 2D reference frames.
    \item Implemented matrix homography and perspective rectification algorithms to correct for camera tilt skew, extracting invariant surface textures from variable viewpoints.
\end{itemize}

\noindent
\textbf{\href{https://satyam-aw.github.io/projects/tweeter/}{Tweeter: Distributed Infrastructure Testbed}} \hfill Sep 2021 – Dec 2021  \\
\textit{High-Concurrency Microblogging Testbed on AWS EC2 Clusters} \hfill UCSB
\begin{itemize}
    \item Conducted high-throughput load testing up to 512 users/sec, mapping user trajectories as state machines to latency scaling curves.
    \item Optimized read throughput by deploying nested fragment caches and database query preloading to eliminate bottleneck points.
\end{itemize}


% ==================== TEACHING ====================
\section*{Teaching}
\noindent
\textbf{University of California, Santa Barbara} \hfill Santa Barbara, CA \\
\textit{\href{https://satyam-aw.github.io/teaching/}{Graduate Teaching Assistant}} \hfill Sep 2021 – Jun 2023
\begin{itemize}
    \item Directed instructional labs, formulated algorithmic programming assignments, and monitored system architecture milestones for senior capstone teams.
    \item Covered advanced and core tracks: CS189 (Capstone), CS184 (Android), CS24 (C++ Data Structures), and CS16 (Problem Solving).
\end{itemize}


% ==================== ACADEMIC SERVICE ====================
\section*{Academic Service \& Peer Review}

\noindent
\textbf{Technical Program Committee \& Peer Reviewer} \\
\textit{ACM Symposium on User Interface Software and Technology (UIST)} \hfill May 2024 \\
Served as an external peer reviewer evaluating full paper submissions for novel interface hardware, interactive system architectures, and spatial tracking methodologies.

\end{document}